\documentclass[11pt]{article}
\usepackage{amsmath, amsthm, amssymb}
\usepackage{mathtools}
\usepackage{hyperref}
\usepackage{geometry}
\geometry{margin=1.2in}

\newtheorem{theorem}{Theorem}[section]
\newtheorem{lemma}[theorem]{Lemma}
\newtheorem{corollary}[theorem]{Corollary}
\newtheorem{definition}[theorem]{Definition}
\theoremstyle{remark}
\newtheorem*{remark}{Remark}

\DeclareMathOperator{\Prob}{\mathbf{Pr}}
\DeclareMathOperator{\Exp}{\mathbf{E}}
\newcommand{\Bool}{\{0,1\}}
\newcommand{\avgFuture}{\mathrm{avgFuture}}
\newcommand{\histTo}{\mathrm{hist}}
\newcommand{\MistakeAt}{\mathrm{MistakeAt}}
\newcommand{\obs}{\mathrm{obs}}
\newcommand{\card}[1]{|#1|}

\title{\textbf{No-Free-Lunch Theorems: Three Variants}\\[0.5em]
\large Natural-Language Proofs}
\author{Proof Agent}
\date{\today}

\begin{document}
\maketitle
\tableofcontents
\newpage

%%% ---------------------------------------------------------------
\section{Overview}

We present natural-language proofs for three variants of the
No-Free-Lunch (NFL) theorem.  All three capture the same meta-theorem
in different settings:

\begin{enumerate}
  \item \textbf{Online adversarial NFL (Theorem~\ref{thm:online})}: no
        deterministic online learner can avoid mistakes against a
        worst-case Boolean target on any pre-fixed query sequence.

  \item \textbf{Wolpert--Macready histogram NFL
        (Theorem~\ref{thm:wm})}: averaged over all cost functions, any
        two non-revisiting stochastic search algorithms see identical
        observation-sequence distributions.

  \item \textbf{Batch stochastic NFL
        (Theorem~\ref{thm:batch})}: for any learning algorithm on a
        finite domain, a realizable distribution exists under which the
        algorithm fails with constant probability.
\end{enumerate}

Throughout, $X$ is the \emph{domain}, $Y$ is the \emph{codomain} (the
range of target / cost functions), and $\Bool$ denotes the set of
Boolean values $\{0,1\}$ (written as \texttt{Bool} in Lean~4).

%%% ---------------------------------------------------------------
\section{Theorem 1: Online Adversarial NFL}
\label{sec:online}

\subsection{Setting}

A \emph{deterministic online learner} $A$ is a function
\[
  A : \mathrm{List}(X \times \Bool) \to X \to \Bool
\]
that receives the history of previously seen (query, label) pairs and
then predicts the label of a new query point.

Given a target $f : X \to \Bool$ and a query sequence
$xs = [x_0, x_1, \ldots, x_{n-1}]$, define the \emph{history at
round $t$}:
\[
  \histTo(f, xs, t) = [(x_0, f(x_0)),\; (x_1, f(x_1)),\; \ldots,\;
    (x_{t-1}, f(x_{t-1}))]
\]
The learner \emph{makes a mistake at round $t$} (written
$\MistakeAt(A, f, xs, t)$) if
\[
  A\!\bigl(\histTo(f, xs, t)\bigr)(x_t) \;\neq\; f(x_t)
  \qquad \text{whenever } t < \card{xs}.
\]

\subsection{Statement}

\begin{theorem}[Adversarial Online NFL]
\label{thm:online}
  For any deterministic online learner $A$ and any
  duplicate-free (\emph{nodup}) query sequence $xs$, there exists a
  target function $f : X \to \Bool$ such that $A$ makes a mistake at
  every round $t$ (i.e., $\MistakeAt(A, f, xs, t)$ holds for all $t$).
\end{theorem}

\subsection{Proof}

By induction on the list $xs$.

\medskip
\noindent\textbf{Base case ($xs = []$).}
The empty sequence has no rounds; any target works vacuously.

\medskip
\noindent\textbf{Inductive step ($xs = x :: xs'$).}
Assume the theorem holds for every algorithm on the tail $xs'$.
Since $xs$ is duplicate-free, $x \notin xs'$ and $xs'$ is itself
duplicate-free.

\medskip
\noindent\emph{Construction.}
\begin{enumerate}
  \item Let $b = \neg A([\,], x)$, the opposite of $A$'s prediction at
        $x$ when it has seen no history.

  \item Define the \emph{shifted learner}
        $A'(h, y) = A\!\bigl((x, b) :: h,\; y\bigr)$.
        Intuitively, $A'$ is $A$ with $(x, b)$ prepended to every
        history.

  \item Apply the induction hypothesis to $A'$ and $xs'$: there exists
        $g : X \to \Bool$ such that $A'$ makes a mistake at every round
        of $xs'$ against $g$.

  \item Define the target:
        \[
          f(y) = \begin{cases} b & \text{if } y = x \\ g(y) &
            \text{otherwise.} \end{cases}
        \]
\end{enumerate}

\medskip
\noindent\emph{Verification that $A$ mistakes at every round of $x :: xs'$
against $f$.}

\medskip
\noindent\textit{Round $t = 0$ (query $x$).}
The history is empty.  $A$'s prediction is $A([\,])(x)$.
By definition $f(x) = b = \neg A([\,])(x)$, so the prediction differs
from the true label.

\medskip
\noindent\textit{Round $t = s+1$ (query $xs'[s]$).}
Since $xs'[s] \in xs'$ and $x \notin xs'$, we have $xs'[s] \neq x$, so
$f(xs'[s]) = g(xs'[s])$.

The history that $A$ sees at round $s+1$ when run on $x :: xs'$ against
$f$ is:
\[
  \bigl[(x, f(x)),\; (xs'[0], f(xs'[0])),\; \ldots,\;
    (xs'[s-1], f(xs'[s-1]))\bigr].
\]
Because $xs'[i] \neq x$ for all $i$, we have $f(xs'[i]) = g(xs'[i])$
for each $i$.  Thus the history equals
\[
  (x, b) ::\; \bigl[(xs'[0], g(xs'[0])),\; \ldots,\;
    (xs'[s-1], g(xs'[s-1]))\bigr]
  = (x, b) ::\; \histTo(g, xs', s).
\]
Therefore:
\[
  A\!\bigl(\histTo(f, x :: xs', s+1)\bigr)(xs'[s])
  = A\!\bigl((x,b) ::\; \histTo(g, xs', s)\bigr)(xs'[s])
  = A'(\histTo(g, xs', s))(xs'[s]).
\]
By the induction hypothesis, $A'$ mistakes at round $s$ of $xs'$ against
$g$, meaning $A'(\histTo(g, xs', s))(xs'[s]) \neq g(xs'[s]) = f(xs'[s])$.
\qed

\medskip
\begin{remark}
The constructed target $f$ is built \emph{adaptively}: it first sets the
label at $x$ to fool $A$ at round 0, then delegates to $g$ (which was
designed to fool the shifted algorithm $A'$ on the remaining queries).
The non-repetition condition ($xs$ is nodup) is essential: if $x$ were
queried again, the consistency of $f$ could not be maintained.
\end{remark}


%%% ---------------------------------------------------------------
\section{Theorem 2: Wolpert--Macready (Histogram Form)}
\label{sec:wm}

\subsection{Setting}

Let $X$ and $Y$ be finite types.  A \emph{search algorithm} $A$ maps a
history $h \in \mathrm{List}(X \times Y)$ to a probability mass function
(PMF) over $X$.  The algorithm is \emph{non-revisiting} if it assigns
zero probability to any point already appearing in the history:
\[
  x \in h.{\rm keys} \;\Rightarrow\; A.\mathrm{nextQuery}(h)(x) = 0.
\]

Given a target cost function $f : X \to Y$, the \emph{$m$-step
observation PMF} of $A$ on $f$ is the distribution over length-$m$
sequences of observed costs $c \in Y^m$ produced when $A$ queries $f$
for $m$ steps (with labels revealed after each query).

\subsection{Statement}

\begin{theorem}[Wolpert--Macready, histogram form]
\label{thm:wm}
  For any two non-revisiting search algorithms $A_1$, $A_2$, finite
  domain $X$, finite codomain $Y$, budget $m \le \card{X}$, and any
  observation sequence $c \in Y^m$:
  \[
    \sum_{f : X \to Y} \Prob\bigl[A_1 \text{ produces } c
      \text{ on } f\bigr]
    \;=\;
    \sum_{f : X \to Y} \Prob\bigl[A_2 \text{ produces } c
      \text{ on } f\bigr].
  \]
\end{theorem}

\subsection{Proof}

\paragraph{Key object.}
For a non-revisiting algorithm $A$, a history $h$ whose query keys are
distinct, and a future observation sequence $c$, define:
\[
  \avgFuture(A, h, c)
  \;=\;
  \sum_{\substack{f : X \to Y \\ f \text{ consistent with } h}}
    \Prob\bigl[A \text{ produces } c \text{ from history } h \text{ on } f\bigr],
\]
where $f$ is \emph{consistent with} $h$ if $f(x) = y$ for every
$(x, y) \in h$.

\begin{lemma}[Algorithm-independent value]
\label{lem:avgfuture}
  If $A$ is non-revisiting, $h$ has distinct query keys, and
  $\card{h} + \card{c} \le \card{X}$, then:
  \[
    \avgFuture(A, h, c) \;=\; \card{Y}^{\card{X} - \card{h} - \card{c}}.
  \]
\end{lemma}

\begin{proof}
By induction on $c$.

\medskip
\noindent\textbf{Base case ($c = [\,]$).}
The empty observation sequence is produced with probability~1 by every
consistent target (there are no steps to take).  Hence:
\[
  \avgFuture(A, h, [\,])
  = \card{\{f : X \to Y \mid f \text{ consistent with } h\}}.
\]
A consistent function must agree with $h$ on the $\card{h}$ queried
points and is free on the remaining $\card{X} - \card{h}$ points.
Therefore the count is $\card{Y}^{\card{X} - \card{h}}$, which equals
$\card{Y}^{\card{X} - \card{h} - 0}$. \checkmark

\medskip
\noindent\textbf{Inductive step ($c = y :: ys$).}
At each step, the algorithm draws the next query $x$ from the PMF
$A.\mathrm{nextQuery}(h)$.  Only if $f(x) = y$ does $x$ contribute to
the observation sequence.  Expanding:
\begin{align*}
  \avgFuture(A, h, y :: ys)
  &= \sum_{\substack{f \text{ consistent} \\ \text{with } h}}
       \sum_{x \in X}
         A.\mathrm{nextQuery}(h)(x)
         \cdot
         \Prob[A \text{ first observes } y \text{ then } ys
               \mid \text{next query is } x, f]
  \\
  &= \sum_{x \in X}
       A.\mathrm{nextQuery}(h)(x)
       \cdot
       \sum_{\substack{f \text{ consistent with } h \\ f(x) = y}}
         \Prob[A \text{ produces } ys \text{ from } h \cup \{(x,y)\} \text{ on } f].
\end{align*}

Two observations justify this rearrangement:
\begin{itemize}
  \item The non-revisiting condition ensures
        $A.\mathrm{nextQuery}(h)(x) = 0$ whenever $x \in h.{\rm keys}$.
        Thus only fresh $x$ contribute.
  \item When $x \notin h.{\rm keys}$, the set of functions consistent
        with $h$ and satisfying $f(x) = y$ is exactly the set of
        functions consistent with $h \cup \{(x, y)\}$.
\end{itemize}

Recognizing the inner double sum as $\avgFuture(A, h \cup \{(x,y)\},
ys)$:
\[
  \avgFuture(A, h, y :: ys)
  = \sum_{x \in X}
      A.\mathrm{nextQuery}(h)(x) \cdot \avgFuture(A, h \cup \{(x,y)\}, ys).
\]
By the induction hypothesis, $\avgFuture(A, h \cup \{(x,y)\}, ys) =
\card{Y}^{\card{X} - (\card{h}+1) - \card{ys}} = \card{Y}^{\card{X} -
\card{h} - \card{c}}$ (a constant independent of $x$).  Since
$A.\mathrm{nextQuery}(h)$ is a probability distribution summing to 1:
\[
  \avgFuture(A, h, y :: ys)
  = \card{Y}^{\card{X} - \card{h} - \card{c}}
    \cdot \underbrace{\sum_{x \in X} A.\mathrm{nextQuery}(h)(x)}_{= 1}
  = \card{Y}^{\card{X} - \card{h} - \card{c}}. \qquad \checkmark
\]
\end{proof}

\paragraph{Conclusion of Theorem~\ref{thm:wm}.}
Apply Lemma~\ref{lem:avgfuture} with $h = [\,]$ (empty history, trivially
consistent with every target) and the given $c$ of length $m \le
\card{X}$.  Every target is consistent with the empty history, so:
\[
  \sum_{f : X \to Y} \Prob[A_i \text{ produces } c \text{ on } f]
  = \avgFuture(A_i, [\,], c)
  = \card{Y}^{\card{X} - m}.
\]
This value is the same for both $A_1$ and $A_2$.
\qed

\begin{corollary}[Averaged performance]
\label{cor:avg}
  Under the same hypotheses, for any score function
  $\varphi : Y^* \to \mathbb{R}_{\ge 0}$:
  \[
    \sum_{f : X \to Y} \Exp[\varphi(A_1 \text{ on } f)]
    \;=\;
    \sum_{f : X \to Y} \Exp[\varphi(A_2 \text{ on } f)].
  \]
\end{corollary}

\begin{proof}
Expand each expectation as a sum over observation sequences $c$ and swap
the two sums:
\begin{align*}
  \sum_{f} \Exp[\varphi(A_i \text{ on } f)]
  &= \sum_{f} \sum_{c \in Y^*}
       \Prob[A_i \text{ produces } c \text{ on } f] \cdot \varphi(c)
  \\
  &= \sum_{c \in Y^*} \varphi(c) \cdot
       \sum_{f} \Prob[A_i \text{ produces } c \text{ on } f].
\end{align*}
By Theorem~\ref{thm:wm}, the inner sum is algorithm-independent, so both
sides are equal.
\qed
\end{proof}

\begin{remark}
The power $\card{Y}^{\card{X} - m}$ is the number of cost functions
consistent with any fixed length-$m$ nodup trace: the values on the $m$
queried points are fixed, and the values on the remaining $\card{X} - m$
points are free.  The proof exploits the fact that in a non-revisiting
run, every trace is automatically nodup.
\end{remark}


%%% ---------------------------------------------------------------
\section{Theorem 3: Batch Stochastic NFL}
\label{sec:batch}

\subsection{Setting}

A \emph{deterministic (batch) learner} $A$ is a function
\[
  A : \mathrm{List}(X \times \Bool) \to (X \to \Bool)
\]
that maps a labelled training sample to a hypothesis.  The
\emph{0--1 loss} of hypothesis $h$ on example $(x, y)$ is $\mathbf{1}[h(x)
\ne y]$.  Given a measure $D$ on $X \times \Bool$, the \emph{risk} of $h$
is $\mathrm{risk}(D, h) = \int \mathbf{1}[h(x) \ne y]\,\mathrm{d}D(x, y)$.

The \emph{failure probability} of $A$ under $D$ with $m$ i.i.d.\ samples
and risk threshold $\varepsilon$ is:
\[
  \mathrm{failureProb}(D, A, m, \varepsilon)
  = \Prob_{S \sim D^m}\!\bigl[\mathrm{risk}(D, A(S)) \ge \varepsilon\bigr].
\]

\subsection{Statement}

\begin{theorem}[Batch stochastic NFL \textup{\cite{ShalevShwartz2014}}]
\label{thm:batch}
  For any deterministic learner $A$ on a finite domain $X$ (with $X$
  nonempty) and any $m$ with $2m \le \card{X}$, there exists a
  distribution $D$ on $X \times \Bool$ such that:
  \begin{enumerate}
    \item \textup{(Realizability)} Some target $f^* : X \to \Bool$
          achieves zero risk under $D$.
    \item \textup{(Hardness)} $\mathrm{failureProb}(D, A, m, 1/8) \ge 1/7$.
  \end{enumerate}
\end{theorem}

\subsection{Proof}

\paragraph{Special case $m = 0$.}
With no training data, the learner $A$ outputs a fixed hypothesis
$h_0 = A([])$.  Pick any $x_0 \in X$ and set $b = \neg h_0(x_0)$.
Let $D = \delta_{(x_0, b)}$ (Dirac mass).
\begin{enumerate}
  \item The constant-$b$ target has zero risk under $D$.
  \item The only sample $S$ of size 0 is the empty tuple; $A([]) = h_0$;
        $\mathrm{risk}(D, h_0) = \mathbf{1}[h_0(x_0) \ne b] = 1 \ge 1/8$.
        Hence $\mathrm{failureProb}(D, A, 0, 1/8) = 1 \ge 1/7$.
\end{enumerate}

\paragraph{Case $m \ge 1$.}
We work in three steps.

\medskip
\noindent\textbf{Step 1: Fixed-sample lower bound on average risk.}

Choose $C \subseteq X$ with $\card{C} = 2m$ (possible since $2m \le
\card{X}$).  Let $Z = \{$elements of $C\}$.

For a labelling $f : Z \to \Bool$ and a sample $S = (z_1, \ldots, z_m)$
of $m$ points from $Z$, define the \emph{abstract risk}:
\[
  R(f, S) = \frac{1}{\card{Z}} \sum_{z \in Z}
    \mathbf{1}\!\bigl[A(\text{sample})(z) \ne f(z)\bigr],
\]
where $\text{sample} = [(z_1, f(z_1)), \ldots, (z_m, f(z_m))]$.

Let $\mathrm{seen}(S) = \{z_1, \ldots, z_m\}$ and
$\mathrm{unseen}(S) = Z \setminus \mathrm{seen}(S)$.  Since
$\card{\mathrm{seen}(S)} \le m$ and $\card{Z} = 2m$, we have
$\card{\mathrm{unseen}(S)} \ge m$.

For any \emph{unseen} point $z \notin \mathrm{seen}(S)$, the label
$f(z)$ does not appear in the training sample.  Therefore the hypothesis
$A(\text{sample})$ is the same regardless of the value of $f(z)$.
Summing the loss at $z$ over both values of $f(z)$ (while holding the
rest of $f$ fixed) gives exactly 1.  Summing over all $2^{\card{Z}-1}$
choices of the other labels:
\[
  \sum_{f : Z \to \Bool} \mathbf{1}\!\bigl[A(\text{sample})(z) \ne f(z)\bigr]
  = 2^{\card{Z}-1} = \frac{2^{\card{Z}}}{2}.
\]

Therefore, for any fixed sample $S$:
\begin{align*}
  \sum_{f : Z \to \Bool} R(f, S)
  &= \frac{1}{\card{Z}} \sum_{z \in Z}
       \sum_{f : Z \to \Bool}
         \mathbf{1}\!\bigl[A(\text{sample})(z) \ne f(z)\bigr]
  \\
  &\ge \frac{1}{\card{Z}} \sum_{z \in \mathrm{unseen}(S)}
       \frac{2^{\card{Z}}}{2}
  \\
  &\ge \frac{\card{\mathrm{unseen}(S)}}{\card{Z}} \cdot \frac{2^{\card{Z}}}{2}
  \\
  &\ge \frac{m}{2m} \cdot \frac{2^{\card{Z}}}{2}
  = \frac{2^{\card{Z}}}{4}.
\end{align*}

\medskip
\noindent\textbf{Step 2: Contradiction argument.}

Suppose for contradiction that every labelling $f : Z \to \Bool$ has
abstract failure rate $< 1/7$, i.e., the fraction of samples $S$ on
which $R(f, S) \ge 1/8$ is less than $1/7$.

For a fixed $f$, let $B_f = \{S : R(f,S) \ge 1/8\}$ be the ``bad''
samples.  Using $R(f, S) \le 1$ always and $R(f, S) < 1/8$ for good
samples:
\begin{align*}
  \sum_{S} R(f, S)
  &\le \sum_{S \in B_f} 1 + \sum_{S \notin B_f} \frac{1}{8}
  \\
  &= \card{B_f} + \frac{\card{\Omega} - \card{B_f}}{8}
  \\
  &= \frac{\card{\Omega}}{8} + \frac{7}{8}\,\card{B_f},
\end{align*}
where $\Omega$ is the set of all samples of size $m$ from $Z$.  If
$\card{B_f}/\card{\Omega} < 1/7$, then $\card{B_f} < \card{\Omega}/7$,
so:
\[
  \sum_{S} R(f, S)
  < \frac{\card{\Omega}}{8} + \frac{7}{8} \cdot \frac{\card{\Omega}}{7}
  = \frac{\card{\Omega}}{8} + \frac{\card{\Omega}}{8}
  = \frac{\card{\Omega}}{4}.
\]

Summing over all $2^{\card{Z}}$ labellings $f$ and swapping the order of
summation:
\[
  \sum_{f} \sum_{S} R(f, S)
  = \sum_{S} \sum_{f} R(f, S)
  \;\ge\; \card{\Omega} \cdot \frac{2^{\card{Z}}}{4}.
\]
But if every $f$ satisfies the failure-rate assumption:
\[
  \sum_{f} \sum_{S} R(f, S) < 2^{\card{Z}} \cdot \frac{\card{\Omega}}{4}
  = \card{\Omega} \cdot \frac{2^{\card{Z}}}{4}.
\]
This is a strict inequality contradicting the lower bound.  Therefore,
some $f^* : Z \to \Bool$ must have abstract failure rate $\ge 1/7$.

\medskip
\noindent\textbf{Step 3: Constructing the hard distribution.}

Extend $f^* : Z \to \Bool$ to a full target $f_X : X \to \Bool$ by
setting $f_X(x) = f^*(x)$ for $x \in C$ and $f_X(x) = 0$ otherwise.
Let $D$ be the uniform distribution on the labelled graph:
\[
  D = \frac{1}{\card{C}} \sum_{x \in C} \delta_{(x,\, f_X(x))}.
\]

\begin{enumerate}
  \item \textbf{Realizability.} The target $f_X$ itself achieves zero
        0--1 loss on every point in the support of $D$, so
        $\mathrm{risk}(D, f_X) = 0$.

  \item \textbf{Failure probability.} The risk of any hypothesis $h$
        under $D$ equals the empirical risk on $C$:
        $\mathrm{risk}(D, h) = \frac{1}{\card{C}} \sum_{x \in C}
        \mathbf{1}[h(x) \ne f_X(x)] = R(f^*, \cdot)$.
        Samples drawn from $D^m$ supported on the graph of $f_X$
        correspond (up to an injective embedding) to support-indexed
        samples.  The failure probability under $D^m$ is therefore at
        least the abstract failure rate $\ge 1/7$.
\end{enumerate}

This completes the proof. \qed

\begin{remark}[Why the constants $1/7$ and $1/8$?]
The argument uses the identity $\frac{7}{8} \cdot \frac{1}{7} =
\frac{1}{8}$.  This ensures that the upper bound
$\frac{\card{\Omega}}{8} + \frac{7}{8}\,\card{B_f}$, evaluated at the
threshold $\card{B_f} = \card{\Omega}/7$, gives exactly $\card{\Omega}/4$
--- matching the lower bound from the unseen-point counting argument.
Any pair $(\varepsilon, \delta)$ with $\varepsilon < 1/2$ and
$(1 - \varepsilon)\,\delta < \varepsilon$ (i.e.\ $\delta < \varepsilon /
(1 - \varepsilon)$) would yield analogous constants.
\end{remark}

\medskip
\begin{thebibliography}{9}
\bibitem{Wolpert1997}
D.~H. Wolpert and W.~G. Macready.
\newblock No free lunch theorems for optimization.
\newblock \emph{IEEE Transactions on Evolutionary Computation}, 1(1):67--82,
  1997.

\bibitem{ShalevShwartz2014}
S.~Shalev-Shwartz and S.~Ben-David.
\newblock \emph{Understanding Machine Learning: From Theory to Algorithms}.
\newblock Cambridge University Press, 2014.
\end{thebibliography}

\end{document}
